\documentclass[aspectratio=169,10pt,spanish]{beamer}

\input{../../../_preambulo_beamer.tex}
\input{../../../_comandos_beamer.tex}

\selectlanguage{spanish}

\title{MA1001B - Modelación Estadística para la Toma de Decisiones}
\subtitle{Lección 06: Técnicas de conteo I}
\author{}
\institute{Tecnol\'ogico de Monterrey}
\date{}

\begin{document}

\begin{frame}[plain]
  \vspace{1.2cm}
  {\Large\bfseries MA1001B - Modelación Estadística para la Toma de Decisiones\par}
  \vspace{0.7cm}
  {\large Lección 06: Técnicas de conteo I\par}
  \vspace{0.7cm}
  {\color{TecRojo}\rule{\textwidth}{1.2pt}}
  \vspace{0.8cm}

  Facultad MA1001B\\[0.3cm]
  Periodo\\[0.3cm]
  Tecnol\'ogico de Monterrey
\end{frame}

\begin{frame}{La pregunta de conteo}
  \begin{alertblock}{Pregunta guía}
    ¿Cómo determinar el tamaño de un espacio muestral finito sin listar cada
    resultado, y cómo detectar el sobreconteo?
  \end{alertblock}

  \vspace{0.5em}
  Al final de esta lección, usted puede:
  \begin{itemize}
    \item aplicar la regla del producto a etapas sucesivas de elección;
    \item distinguir asignaciones ordenadas de subconjuntos no ordenados;
    \item verificar fórmulas mediante enumeración exhaustiva;
    \item diagnosticar conteo repetido cuando las restricciones interactúan.
  \end{itemize}

  \vfill
  \scriptsize Todas las instalaciones, roles, personas y restricciones son sintéticos.
\end{frame}

\begin{frame}{Paso 1: Regla del producto para configuraciones completas}
  \begin{columns}[T]
    \begin{column}{0.40\textwidth}
      \small
      Elegir un elemento de cada etapa:
      \[
        N=\prod_{i=1}^{m}n_i.
      \]

      \begin{block}{Configuración sintética de revisión}
        3 centros de cumplimiento\\
        4 etapas de flujo de trabajo\\
        2 tipos de evidencia
      \end{block}

      \[
        3\times4\times2=24
      \]

      La enumeración exhaustiva en Python también devuelve 24 ternas. Cada
      centro aporta 8 configuraciones.
    \end{column}
    \begin{column}{0.60\textwidth}
      \centering
      \includegraphics[width=\textwidth]{../figuras/conteo_01_regla_producto.png}
    \end{column}
  \end{columns}
\end{frame}

\begin{frame}{Paso 2: Permutaciones para roles distintos}
  \begin{columns}[T]
    \begin{column}{0.40\textwidth}
      \small
      Asignar roles de Líder, Validación y Reportes a partir de seis
      analistas.

      \[
        P(n,k)=\frac{n!}{(n-k)!}
      \]
      \[
        P(6,3)=6\times5\times4=120
      \]

      \begin{exampleblock}{Por qué importa el orden}
        Los mismos tres analistas en roles distintos crean una asignación
        diferente. Python genera exactamente 120 ternas ordenadas.
      \end{exampleblock}
    \end{column}
    \begin{column}{0.60\textwidth}
      \centering
      \includegraphics[width=\textwidth]{../figuras/conteo_02_permutaciones.png}
    \end{column}
  \end{columns}
\end{frame}

\begin{frame}{Paso 3: Combinaciones para un panel sin ranking}
  \begin{columns}[T]
    \begin{column}{0.40\textwidth}
      \small
      Seleccionar tres analistas para un panel de pares sin roles internos.

      \[
        \binom{n}{k}=\frac{n!}{k!(n-k)!}
      \]
      \[
        \binom{6}{3}=20
      \]

      Cada subconjunto de tres personas tiene $3!=6$ ordenamientos posibles:
      \[
        \frac{P(6,3)}{3!}=\frac{120}{6}=20.
      \]
    \end{column}
    \begin{column}{0.60\textwidth}
      \centering
      \includegraphics[width=\textwidth]{../figuras/conteo_03_combinaciones.png}
    \end{column}
  \end{columns}
\end{frame}

\begin{frame}{El resultado determina el método de conteo}
  \small
  \begin{center}
    \begin{tabular}{p{0.27\textwidth}p{0.26\textwidth}p{0.34\textwidth}}
      \toprule
      \textbf{Pregunta} & \textbf{Método} & \textbf{Resultado de la Lección 06} \\
      \midrule
      Una elección en cada etapa & Regla del producto & $3\times4\times2=24$ configuraciones \\
      Posiciones distintas asignadas & Permutación & $P(6,3)=120$ asignaciones de roles \\
      Subconjunto sin ranking & Combinación & $\binom{6}{3}=20$ paneles \\
      \bottomrule
    \end{tabular}
  \end{center}

  \vspace{0.8em}
  \begin{alertblock}{Prueba de decisión}
    Describa un resultado final antes de elegir una fórmula. Si intercambiar
    dos elementos seleccionados cambia ese resultado, el orden importa.
  \end{alertblock}

  Una fórmula puede devolver un entero y aun así responder la pregunta de
  conteo equivocada.
\end{frame}

\begin{frame}{Paso 4: El límite del sobreconteo}
  \begin{columns}[T]
    \begin{column}{0.43\textwidth}
      \small
      Elegir cuatro personas de cinco analistas de operaciones y cuatro
      especialistas de datos. Exigir ambas disciplinas y excluir el par O1
      con D1.

      \vspace{0.4em}
      \begin{alertblock}{Expresión ingenua}
        $5\times4\times\binom{7}{2}=420$\\
        El universo completo tiene solo $\binom{9}{4}=126$ grupos.
      \end{alertblock}

      \[
        126-5-1-21=99
      \]
      La enumeración exhaustiva confirma 99 grupos válidos.
    \end{column}
    \begin{column}{0.57\textwidth}
      \centering
      \includegraphics[width=\textwidth]{../figuras/conteo_04_restricciones.png}
    \end{column}
  \end{columns}
\end{frame}

\begin{frame}{Verificación de conceptos}
  \begin{enumerate}
    \item Un código de revisión usa 4 centros, 3 etapas y 2 tipos de
      evidencia. ¿Cuántas configuraciones completas existen?
    \item ¿Cuántas asignaciones ordenadas llenan dos roles distintos a
      partir de cinco personas?
    \item ¿Cuántos paneles de dos personas sin ranking se pueden seleccionar
      de cinco personas?
    \item ¿Por qué un resultado de 420 revela de inmediato un error cuando el
      universo completo contiene solo 126 grupos únicos?
  \end{enumerate}

  \vspace{0.5em}
  \begin{block}{Conexión con la Lección 07}
    La siguiente lección extiende el conteo a categorías repetidas y
    asignaciones multinomiales, donde varios elementos comparten el mismo
    tipo.
  \end{block}

  \vfill
  \scriptsize Fuente: OpenStax, \textit{Introductory Business Statistics 2e},
  Capítulo 3, Sección 3.1. CC BY-NC-SA 4.0.
\end{frame}

\appendix



\end{document}
